\chapter{Mecanismos de reação: as setas curvas}\label{ch:g12:curly-arrows}

Uma \omterm{def:g8:balanced-equations:equation}{equação balanceada} é como a primeira e a última fotografia de um
filme: ela mostra os \omterm{def:g7:chemical-reactions:reactant}{reagentes} antes e os produtos depois, e nada no
meio. No entanto, as \omterm{def:g7:atoms-and-molecules:molecule}{moléculas} não se limitam a trocar \omterm{def:g7:atoms-and-molecules:atom}{átomos}. As
ligações se rompem e se formam uma de cada vez, à medida que pares de
\omterm{def:g9:inside-the-atom:nucleus}{elétrons} se deslocam de um lugar para outro, passando por espécies de
vida curta que nunca aparecem na equação. Os químicos desenham esses
movimentos com \omterm{def:g12:curly-arrows:curly-arrow}{setas curvas}. Com algumas regras, essas setas explicam
por que uma reação acontece onde acontece, e preveem reações nunca vistas
antes.

\begin{recall}
A \omterm{def:g11:polarity-and-cohesion:electronegativity}{eletronegatividade}, as \omterm{def:g11:polarity-and-cohesion:polar-bond}{ligações polares} e as \omterm{def:g11:polarity-and-cohesion:polar-bond}{cargas parciais}
(\cref{ch:g11:polarity-and-cohesion}). As \omterm{def:g10:lewis-and-shape:lewis-structure}{estruturas de Lewis}: \omterm{def:g10:lewis-and-shape:lone-pair}{pares
ligantes} e \omterm{def:g10:lewis-and-shape:lone-pair}{pares não ligantes} (\cref{ch:g10:lewis-and-shape}). Os \omterm{def:g11:functional-groups:functional-group}{grupos
funcionais} e as funções (\cref{ch:g11:functional-groups}).
\end{recall}

\section{Sítios ricos e sítios pobres em elétrons}

\begin{definition}[Sítios doadores e sítios aceptores de elétrons]\label{def:g12:curly-arrows:site}
Um \emph{sítio doador de elétrons}\index{sitio doador de eletrons@sítio doador de elétrons} de uma
\omterm{def:g7:atoms-and-molecules:molecule}{molécula} ou de um \omterm{def:g9:ions:ion}{íon} é um lugar rico em \omterm{def:g9:inside-the-atom:nucleus}{elétrons}, capaz de ceder um par
para formar uma nova ligação: um \omterm{def:g10:lewis-and-shape:lone-pair}{par não ligante}, uma \omterm{def:g10:lewis-and-shape:multiple-bond}{ligação dupla}, um
\omterm{def:g7:atoms-and-molecules:atom}{átomo} com carga negativa ou com \omterm{def:g11:polarity-and-cohesion:polar-bond}{carga parcial} $\delta^-$. Um \emph{sítio
aceptor de elétrons}\index{sitio aceptor de eletrons@sítio aceptor de elétrons} é um lugar pobre em
\omterm{def:g9:inside-the-atom:nucleus}{elétrons}, capaz de receber um par: um \omterm{def:g7:atoms-and-molecules:atom}{átomo} com carga positiva ou com
\omterm{def:g11:polarity-and-cohesion:polar-bond}{carga parcial} $\delta^+$.
\end{definition}

\begin{example}[Os sítios de duas moléculas]\label{ex:g12:curly-arrows:sites}
No bromoetano, \ce{CH3-CH2-Br}, o bromo é mais eletronegativo que o
carbono: o carbono ligado a ele tem $\delta^+$ (sítio aceptor), o bromo
$\delta^-$ e três \omterm{def:g10:lewis-and-shape:lone-pair}{pares não ligantes}. Na propanona, o oxigênio do grupo
\ce{C=O} tem $\delta^-$ e dois \omterm{def:g10:lewis-and-shape:lone-pair}{pares não ligantes} (sítios doadores), o
seu carbono $\delta^+$ (sítio aceptor). O \omterm{def:g9:ions:ion}{íon} hidróxido \ce{OH-}, com a
sua carga negativa e os seus três \omterm{def:g10:lewis-and-shape:lone-pair}{pares não ligantes}, é um doador forte.
\end{example}

\begin{omfigure}
\setchemfig{atom sep=2.2em}
\begin{tabular}{cc}
\chemfig{H_3C-C(-[2]H)(-[6]H)-\charge{90:2pt=\:,0=\:,270:2pt=\:}{Br}} &
\chemfig{H_3C-C(-[6]CH_3)=[2]\charge{0=\:,180=\:}{O}} \\[2pt]
\small $\delta^+$ no carbono ligado ao \ce{Br}, $\delta^-$ no \ce{Br} &
\small $\delta^+$ no carbono do \ce{C=O}, $\delta^-$ no \ce{O} \\
\small bromoetano & \small propanona
\end{tabular}

{\small \omterm{def:g7:atoms-and-molecules:atom}{Átomos} de carbono pobres em \omterm{def:g9:inside-the-atom:nucleus}{elétrons} ($\delta^+$, sítios
aceptores) ao lado de \omterm{def:g7:atoms-and-molecules:atom}{átomos} eletronegativos ricos em \omterm{def:g9:inside-the-atom:nucleus}{elétrons}
($\delta^-$, com \omterm{def:g10:lewis-and-shape:lone-pair}{pares não ligantes}: sítios doadores).}
\end{omfigure}

\section{As setas curvas}

\begin{definition}[Seta curva]\label{def:g12:curly-arrows:curly-arrow}
Uma \emph{seta curva}\index{seta curva} mostra o deslocamento de um par
de \omterm{def:g9:inside-the-atom:nucleus}{elétrons} durante uma etapa de uma reação. Ela parte do par que se
desloca, um \omterm{def:g10:lewis-and-shape:lone-pair}{par não ligante} ou uma ligação, e termina onde o par vai: num
\omterm{def:g7:atoms-and-molecules:atom}{átomo}, onde ele forma uma nova ligação, ou num \omterm{def:g7:atoms-and-molecules:atom}{átomo} que fica com o par
de uma ligação que se rompe.
\end{definition}

\begin{method}[Desenhar as setas curvas]\label{met:g12:curly-arrows:drawing}
\begin{enumerate}
  \item Encontre o sítio doador e o sítio aceptor.
  \item Desenhe uma seta do par doador (um \omterm{def:g10:lewis-and-shape:lone-pair}{par não ligante} ou uma
        ligação) até o \omterm{def:g7:atoms-and-molecules:atom}{átomo} aceptor: ali se forma uma nova ligação.
  \item Se o \omterm{def:g7:atoms-and-molecules:atom}{átomo} aceptor ficar então com \omterm{def:g9:inside-the-atom:nucleus}{elétrons} demais (mais que um
        octeto, ou mais que dois para o hidrogênio), desenhe uma segunda
        seta que rompe uma das suas ligações, com o par indo para o \omterm{def:g7:atoms-and-molecules:atom}{átomo}
        mais eletronegativo.
  \item Confira as cargas depois da etapa: o \omterm{def:g7:atoms-and-molecules:atom}{átomo} que cedeu um \omterm{def:g10:lewis-and-shape:lone-pair}{par não
        ligante} fica uma unidade mais positivo, o \omterm{def:g7:atoms-and-molecules:atom}{átomo} que recebeu um
        par como \omterm{def:g10:lewis-and-shape:lone-pair}{par não ligante}, uma unidade mais negativo; a carga
        total não muda.
\end{enumerate}
\end{method}

\begin{proposition}[A carga se conserva em cada etapa]\label{prop:g12:curly-arrows:charge}
A carga elétrica total das espécies é a mesma antes e depois de cada
etapa de um mecanismo.
\end{proposition}

\begin{proof}
Uma \omterm{def:g12:curly-arrows:curly-arrow}{seta curva} só desloca \omterm{def:g9:inside-the-atom:nucleus}{elétrons} que já estão ali; nenhum \omterm{def:g9:inside-the-atom:nucleus}{elétron} é
criado nem destruído, e os \omterm{def:g9:inside-the-atom:nucleus}{núcleos} não mudam.
\end{proof}

\begin{example}[Uma primeira seta]\label{ex:g12:curly-arrows:first}
Um \omterm{def:g9:ions:ion}{íon} hidrogênio \ce{H+} encontra um \omterm{def:g9:ions:ion}{íon} hidróxido \ce{OH-}. Um \omterm{def:g10:lewis-and-shape:lone-pair}{par não
ligante} do oxigênio forma uma ligação com o \omterm{def:g9:ions:ion}{íon} hidrogênio, que não tem
nenhum \omterm{def:g9:inside-the-atom:nucleus}{elétron}: \ce{H+ + OH- -> H2O}. Uma seta, do \omterm{def:g10:lewis-and-shape:lone-pair}{par não ligante} do
oxigênio até o \ce{H+}; as cargas $+1$ e $-1$ viram 0 na \omterm{def:g7:atoms-and-molecules:molecule}{molécula} de
água.
\end{example}

\begin{omfigure}
\setchemfig{atom sep=2.1em}
\chemfig{@{hp}H^{+}}\qquad\raisebox{0.2ex}{$+$}\qquad
\chemfig{H-@{ox}\charge{90=\:,0=\:,270=\:}{O}^{-}}
\qquad\raisebox{0.2ex}{$\longrightarrow$}\qquad
\chemfig{H-\charge{90=\:,270=\:}{O}-H}
\chemmove{
  \draw[->, thick, omProp, shorten <=4pt, shorten >=3pt]
    ($(ox)+(0,0.25)$) .. controls +(110:9mm) and +(60:9mm) .. (hp);}

{\small A etapa mais simples: um \omterm{def:g10:lewis-and-shape:lone-pair}{par não ligante} do \omterm{def:g9:ions:ion}{íon} hidróxido se
desloca para formar a nova ligação \ce{O-H} com o \omterm{def:g9:ions:ion}{íon} hidrogênio.}
\end{omfigure}

\section{Etapas elementares e intermediários}

\begin{definition}[Mecanismo, etapa elementar, intermediário]\label{def:g12:curly-arrows:mechanism}
O \emph{mecanismo de reação}\index{mecanismo de reacao@mecanismo de reação} de uma reação é
a série de etapas pelas quais os \omterm{def:g7:chemical-reactions:reactant}{reagentes} se tornam os produtos. Cada
\emph{etapa elementar}\index{etapa elementar} é um único evento em que
alguns pares de \omterm{def:g9:inside-the-atom:nucleus}{elétrons} se deslocam ao mesmo tempo. Uma espécie formada
numa etapa e consumida numa etapa posterior é um \emph{intermediário de
reação}\index{intermediario de reacao@intermediário de reação}: ela não aparece na equação
global.
\end{definition}


\begin{example}[Um carbocátion como intermediário]\label{ex:g12:curly-arrows:carbocation}
Quando o brometo de hidrogênio se adiciona ao eteno, a reação tem duas
etapas. Na primeira, a \omterm{def:g10:lewis-and-shape:multiple-bond}{ligação dupla} do eteno cede um par ao hidrogênio
do \ce{HBr}, e o par \ce{H-Br} vai para o bromo, que sai como \ce{Br-};
o carbono que perdeu a sua parte da \omterm{def:g10:lewis-and-shape:multiple-bond}{ligação dupla} fica com só seis
\omterm{def:g9:inside-the-atom:nucleus}{elétrons} e uma carga positiva: um carbocátion, \ce{CH3-CH2+}. Na segunda
etapa, um \omterm{def:g10:lewis-and-shape:lone-pair}{par não ligante} do \ce{Br-} forma uma ligação com esse
carbono. O carbocátion é um intermediário.
\end{example}

\section{Três categorias de reação}

\begin{definition}[Substituição, adição, eliminação]\label{def:g12:curly-arrows:categories}
Numa \emph{substituição}\index{substituicao@substituição}, um \omterm{def:g7:atoms-and-molecules:atom}{átomo} ou um grupo de
uma \omterm{def:g7:atoms-and-molecules:molecule}{molécula} é trocado por outro. Numa \emph{adição}\index{adicao@adição}, duas
\omterm{def:g7:atoms-and-molecules:molecule}{moléculas} se juntam numa só, e uma ligação múltipla vira simples. Numa
\emph{eliminação}\index{eliminacao@eliminação}, uma \omterm{def:g7:atoms-and-molecules:molecule}{molécula} pequena (muitas vezes a
água) é retirada de uma \omterm{def:g7:atoms-and-molecules:molecule}{molécula}, e uma ligação simples vira múltipla.
\end{definition}

\begin{omfigure}
\setchemfig{atom sep=2em}
\small
\textbf{\omterm{def:g12:curly-arrows:categories}{Substituição}} (uma etapa): \ce{CH3-CH2-Br + OH- -> CH3-CH2-OH + Br-}

\medskip
\chemfig{H-@{o1}\charge{90=\:,0=\:,270=\:}{O}^{-}}\qquad
\chemfig{H_3C-@{c1}C(-[2]H)(-[6]H)-[@{b1}]@{br1}\charge{90=\:,0=\:,270=\:}{Br}}
\qquad$\longrightarrow$\qquad
\chemfig{H_3C-C(-[2]H)(-[6]H)-O-H}\quad$+$\quad\ce{Br-}
\chemmove{
  \draw[->, thick, omProp, shorten <=4pt, shorten >=3pt]
    ($(o1)+(0.15,0.2)$) .. controls +(60:9mm) and +(120:9mm) .. (c1);
  \draw[->, thick, omDef, shorten <=2pt, shorten >=4pt]
    (b1) .. controls +(-90:7mm) and +(-90:7mm) .. (br1);}

\vspace{6mm}
\textbf{\omterm{def:g12:curly-arrows:categories}{Adição}} (duas etapas): \ce{CH2=CH2 + HBr -> CH3-CH2-Br}

\vspace{9mm}
\chemfig{H_2C=[@{pi}]CH_2}\quad$+$\quad
\chemfig{@{h2}H-[@{b2}]@{br2}Br}
\qquad$\longrightarrow$\qquad
\chemfig{H_3C-@{cp}\charge{90:2pt=$\scriptstyle +$}{C}H_2}\quad$+$\quad
\chemfig{@{brm}\charge{90=\:,180=\:,270=\:,0=\:}{Br}^{-}}
\qquad$\longrightarrow$\qquad \ce{CH3-CH2-Br}
\chemmove{
  \draw[->, thick, omProp, shorten <=2pt, shorten >=3pt]
    (pi) .. controls +(90:10mm) and +(110:9mm) .. (h2);
  \draw[->, thick, omDef, shorten <=2pt, shorten >=4pt]
    (b2) .. controls +(-90:7mm) and +(-90:7mm) .. (br2);
  \draw[->, thick, omProp, shorten <=4pt, shorten >=3pt]
    ($(brm)+(0,-0.25)$) .. controls +(-120:8mm) and +(-60:8mm) .. (cp);}

\vspace{6mm}
\textbf{\omterm{def:g12:curly-arrows:categories}{Eliminação}} (desidratação ácida do etanol, três etapas):
\ce{CH3-CH2-OH -> CH2=CH2 + H2O}

\medskip
\begin{tabular}{@{}l@{}}
1.\ o oxigênio capta um \omterm{def:g9:ions:ion}{íon} hidrogênio: \ce{CH3-CH2-OH + H+ -> CH3-CH2-OH2+};\\
2.\ o par \ce{C-O} sai com a água: \ce{CH3-CH2-OH2+ -> CH3-CH2+ + H2O};\\
3.\ o par de uma ligação \ce{C-H} vizinha vira a \omterm{def:g10:lewis-and-shape:multiple-bond}{ligação dupla}, e o
    \ce{H+} sai:\\
\phantom{3.\ }\ce{CH3-CH2+ -> CH2=CH2 + H+}.
\end{tabular}

\medskip
{\small Três mecanismos. Setas laranja: um par doador forma uma nova
ligação. Setas azuis: uma ligação se rompe, e o seu par vai para o \omterm{def:g7:atoms-and-molecules:atom}{átomo}
mais eletronegativo. Na \omterm{def:g12:curly-arrows:categories}{eliminação}, o \omterm{def:g9:ions:ion}{íon} hidrogênio é devolvido no
final.}
\end{omfigure}


\begin{remark}[Mudar o esqueleto ou o grupo]\label{rem:g12:curly-arrows:skeleton}
Algumas reações só mudam o \omterm{def:g11:functional-groups:functional-group}{grupo funcional} e conservam o esqueleto
carbônico: a \omterm{def:g12:curly-arrows:categories}{substituição} do bromoetano transforma um \omterm{def:g11:functional-groups:halogenoalkane}{haloalcano} num
\omterm{def:g11:functional-groups:alcohol}{álcool}, dois carbonos antes, dois depois. Outras constroem ou cortam o
esqueleto: formar uma nova ligação carbono--carbono alonga uma cadeia, a
etapa-chave para fabricar uma \omterm{def:g7:atoms-and-molecules:molecule}{molécula} grande a partir de \omterm{def:g7:atoms-and-molecules:molecule}{moléculas}
pequenas. Planejar uma \omterm{def:g10:synthesis-yield:synthesis}{síntese} é escolher reações dos dois tipos, na
ordem certa.
\end{remark}

\section{Exercícios}

\begin{exercise}[$\star$]\label{exo:g12:curly-arrows:1}
Em cada espécie, aponte um sítio doador e, se houver, um sítio aceptor:
\ce{OH-}; \ce{H+}; \ce{CH2=CH2}; \ce{CH3-Cl}; \ce{H2O}.
\end{exercise}

\begin{exercise}[$\star$]\label{exo:g12:curly-arrows:2}
\omterm{def:g12:curly-arrows:categories}{Substituição}, \omterm{def:g12:curly-arrows:categories}{adição} ou \omterm{def:g12:curly-arrows:categories}{eliminação}?
(a) \ce{CH3-CH2-Cl + OH- -> CH3-CH2-OH + Cl-};
(b) \ce{CH2=CH2 + H2O -> CH3-CH2-OH};
(c) \ce{CH3-CH2-OH -> CH2=CH2 + H2O};
(d) \ce{CH2=CH2 + Br2 -> CH2Br-CH2Br}.
\end{exercise}

\begin{exercise}[$\star$]\label{exo:g12:curly-arrows:3}
De onde parte uma \omterm{def:g12:curly-arrows:curly-arrow}{seta curva}, e onde ela termina? O que significa uma
\omterm{def:g12:curly-arrows:curly-arrow}{seta curva} que vai de um \omterm{def:g10:lewis-and-shape:lone-pair}{par não ligante} de um \omterm{def:g7:atoms-and-molecules:atom}{átomo} de oxigênio até um
\omterm{def:g7:atoms-and-molecules:atom}{átomo} de carbono?
\end{exercise}

\begin{exercise}[$\star$]\label{exo:g12:curly-arrows:4}
Na reação \ce{H+ + NH3 -> NH4+}, que espécie cede o par? Desenhe a seta.
Qual é a carga do produto?
\end{exercise}

\begin{exercise}[$\star$]\label{exo:g12:curly-arrows:5}
O que é um \omterm{def:g12:curly-arrows:mechanism}{intermediário de reação}? Dê o nome do intermediário da \omterm{def:g12:curly-arrows:categories}{adição}
do \ce{HBr} ao eteno.
\end{exercise}

\begin{exercise}[$\star\star$]\label{exo:g12:curly-arrows:6}
Desenhe as \omterm{def:g12:curly-arrows:curly-arrow}{setas curvas} da \omterm{def:g12:curly-arrows:categories}{substituição}
\ce{CH3-Br + OH- -> CH3-OH + Br-} (uma etapa), e confira as cargas.
\end{exercise}

\begin{exercise}[$\star\star$]\label{exo:g12:curly-arrows:7}
Na primeira etapa da \omterm{def:g12:curly-arrows:categories}{adição} do \ce{HBr} ao eteno, por que o par
\ce{H-Br} vai para o bromo e não para o hidrogênio? Que carga tem cada
produto da etapa?
\end{exercise}

\begin{exercise}[$\star\star$]\label{exo:g12:curly-arrows:8}
Na desidratação do etanol, escreva as três etapas e diga, para cada uma,
que espécie é o doador e qual é o aceptor. Por que o \omterm{def:g9:ions:ion}{íon} hidrogênio não
é escrito na equação global?
\end{exercise}

\begin{exercise}[$\star\star$]\label{exo:g12:curly-arrows:9}
A etapa \ce{CH3-CH2+ + H2O -> CH3-CH2-OH2+} vem depois da \omterm{def:g12:curly-arrows:categories}{adição} de um
\omterm{def:g9:ions:ion}{íon} hidrogênio ao eteno em água. Desenhe a seta. O que se forma depois
de mais uma perda de \ce{H+}? Escreva a equação global e dê a sua
categoria.
\end{exercise}

\begin{exercise}[$\star\star$]\label{exo:g12:curly-arrows:10}
Leia o mecanismo da \omterm{def:g12:curly-arrows:categories}{substituição} do bromoetano: quantos pares se
deslocam? Que ligação se forma, qual se rompe? Há um intermediário?
\end{exercise}

\begin{exercise}[$\star\star$]\label{exo:g12:curly-arrows:11}
% ledger: en:C, en:Cl, en:Br
O clorometano reage com o \omterm{def:g9:ions:ion}{íon} hidróxido mais devagar que o bromometano.
Usando as \omterm{def:g11:polarity-and-cohesion:electronegativity}{eletronegatividades} (C 2.55, Cl 3.16, Br 2.96), explique qual
carbono é o mais $\delta^+$. Isso explica as velocidades? (Pense também
na facilidade com que o \omterm{def:g9:ions:ion}{íon} que sai leva o par.)
\end{exercise}

\begin{exercise}[$\star\star\star$]\label{exo:g12:curly-arrows:12}
O cloreto de hidrogênio se adiciona ao propeno, \ce{CH2=CH-CH3}. Na
primeira etapa, o \omterm{def:g9:ions:ion}{íon} hidrogênio pode se ligar a qualquer um dos
carbonos da \omterm{def:g10:lewis-and-shape:multiple-bond}{ligação dupla}, o que dá dois carbocátions possíveis. Desenhe
os dois, os dois produtos possíveis, e dê o nome deles. (Qual se forma
mais é estudado no volume do primeiro ano da universidade.)
\end{exercise}

\begin{exercise}[$\star\star\star$]\label{exo:g12:curly-arrows:13}
Um aluno desenha a primeira etapa da \omterm{def:g12:curly-arrows:categories}{substituição} do bromoetano com uma
seta que vai do \omterm{def:g7:atoms-and-molecules:atom}{átomo} de bromo até o oxigênio do \ce{OH-}. Explique os
dois erros.
\end{exercise}

\begin{exercise}[$\star\star\star$]\label{exo:g12:curly-arrows:14}
Na formação de um \omterm{def:g11:functional-groups:carboxylic-acid}{éster} a partir de um \omterm{def:g11:functional-groups:carboxylic-acid}{ácido carboxílico} e de um \omterm{def:g11:functional-groups:alcohol}{álcool},
o oxigênio do \omterm{def:g11:functional-groups:alcohol}{álcool} ataca o carbono do grupo \ce{C=O}. Usando as \omterm{def:g11:polarity-and-cohesion:polar-bond}{cargas
parciais}, explique por que esse carbono é um sítio aceptor e esse
oxigênio um sítio doador. Que \omterm{def:g7:atoms-and-molecules:molecule}{molécula} pequena é perdida no total, e de
que categoria é a reação inteira?
\end{exercise}

\begin{exercise}[$\star\star\star$]\label{exo:g12:curly-arrows:15}
O eteno reage com o bromo \ce{Br2}, uma \omterm{def:g11:polarity-and-cohesion:polar-molecule}{molécula apolar}. Quando uma
\omterm{def:g7:atoms-and-molecules:molecule}{molécula} de \ce{Br2} se aproxima da \omterm{def:g10:lewis-and-shape:multiple-bond}{ligação dupla}, os \omterm{def:g9:inside-the-atom:nucleus}{elétrons} da
ligação \ce{Br-Br} são empurrados para longe do \omterm{def:g7:atoms-and-molecules:atom}{átomo} de bromo mais
próximo. Explique por que isso cria um sítio aceptor, e proponha a
primeira etapa do mecanismo.
\end{exercise}

\section{Problema: fabricar o bromoetano}

\begin{problem}[{Problema de fim de semana --- a partir de dez gramas de
eteno: quanto bromoetano?}]
\label{pb:g12:curly-arrows:1}
% ledger: aw:C, aw:H, aw:Br, en:H, en:Br, en:C
O bromoetano, material de partida de muitas \omterm{def:g10:synthesis-yield:synthesis}{sínteses}, é fabricado pela
\omterm{def:g12:curly-arrows:categories}{adição} de brometo de hidrogênio ao eteno:
\ce{CH2=CH2 + HBr -> CH3-CH2-Br}. Um químico usa \qty{10.0}{g} de eteno
e um excesso de brometo de hidrogênio; o \omterm{def:g10:synthesis-yield:yield}{rendimento} é de
\qty{75}{\percent}.

\medskip
\noindent\textbf{Parte I --- Os sítios.}
\begin{enumerate}
  \item Desenhe as \omterm{def:g10:lewis-and-shape:lewis-structure}{estruturas de Lewis} do eteno e do brometo de
        hidrogênio.
  \item Onde fica o sítio doador do eteno? Por quê?
  \item Calcule a diferença de \omterm{def:g11:polarity-and-cohesion:electronegativity}{eletronegatividade} da ligação \ce{H-Br}
        (H 2.20, Br 2.96). Que \omterm{def:g7:atoms-and-molecules:atom}{átomo} tem $\delta^+$?
  \item Que \omterm{def:g7:atoms-and-molecules:atom}{átomo} do \ce{HBr} é o sítio aceptor?
\end{enumerate}

\medskip
\noindent\textbf{Parte II --- O mecanismo.}
\begin{enumerate}[resume]
  \item Desenhe a primeira etapa com as suas duas \omterm{def:g12:curly-arrows:curly-arrow}{setas curvas}.
  \item Dê as duas espécies formadas e as suas cargas.
  \item Verifique a conservação da carga na primeira etapa.
  \item Desenhe a segunda etapa com a sua \omterm{def:g12:curly-arrows:curly-arrow}{seta curva}.
  \item Quantos \omterm{def:g9:inside-the-atom:nucleus}{elétrons} cercam o carbono positivo do intermediário? Por
        que ele é tão reativo?
\end{enumerate}

\medskip
\noindent\textbf{Parte III --- Que tipo de reação?}
\begin{enumerate}[resume]
  \item Qual é o intermediário do mecanismo?
  \item Por que ele não aparece na equação global?
  \item A que categoria pertence a reação?
  \item A reação muda o esqueleto carbônico, ou só o \omterm{def:g11:functional-groups:functional-group}{grupo funcional}?
\end{enumerate}

\medskip
\noindent\textbf{Parte IV --- Quanto produto?}
\begin{enumerate}[resume]
  \item Calcule as \omterm{def:g10:the-mole:molar-mass}{massas molares} do eteno e do bromoetano.
  \item Calcule a \omterm{def:g10:the-mole:amount}{quantidade de matéria} de eteno.
  \item Preencha a \omterm{def:g10:reaction-progress-table:extent}{tabela de avanço}. Qual \omterm{def:g7:chemical-reactions:reactant}{reagente} é o limitante?
  \item Calcule a massa máxima de bromoetano.
  \item Calcule a massa obtida com um \omterm{def:g10:synthesis-yield:yield}{rendimento} de \qty{75}{\percent}.
  \item Dê a resposta final: que massa de bromoetano o químico obtém?
\end{enumerate}
\end{problem}
